% Jackiw-Teitelboim Gravity from Counting Causally Convex Subsets
% Target: Physical Review D
% Author: Thomas DiFiore

\documentclass[aps,prd,twocolumn,superscriptaddress]{revtex4-2}

\usepackage{amsmath,amssymb,amsfonts}
\usepackage{graphicx}
\usepackage{booktabs}
\usepackage{hyperref}

\begin{document}

\title{Jackiw-Teitelboim Gravity and Hagedorn Density of States\\from Counting Causally Convex Subsets}

\author{Thomas DiFiore}

\date{\today}

\begin{abstract}
We define a partition function over causally convex subsets of a
discrete spacetime, weighted by the Benincasa-Dowker action:
$Z(\beta) = \sum_{S \in \mathrm{CC}(P)} e^{-\beta \cdot S_{\mathrm{BD}}[S]}$.
On 42 variable-width two-dimensional grids at fixed element count
$n = 15$ and curvatures ranging from $R = -10$ to $R = +10$, we find
that the effective action decomposes as
$S_{\mathrm{eff}} = a \cdot R + b \cdot \sum (\log w_{i+1} -
\log w_i)^2$. A massive-scale test (2{,}061 profiles across
$n = 15$ and $n = 20$) confirms the decomposition with
$R^2 \approx 0.71$ at $\beta = 1$ rising to $\approx 0.81$ at
$\beta = 3$, with the log-Schwarzian dominating
($R^2 \approx 0.58$ alone) and the curvature $R$ contributing only
$\sim 10\%$. At $n = 30$ (50 profiles), the topological
Einstein-Hilbert term $R$ becomes negligible ($R^2 < 0.01$) while
the conformal-mode kinetic term (the log-Schwarzian) explains
$65\%$ of the variance alone. Both coefficients scale linearly with
$\beta$ ($R^2 > 0.99$) at both system sizes, supporting the
decomposition into
$\log Z = -\beta \cdot S_{\mathrm{eff}} + \mathrm{entropy}$.
The vanishing of the $R$ contribution at larger $n$ is the expected
2D behavior: the Einstein-Hilbert action is topological
(Gauss-Bonnet) and non-dynamical. The log-Schwarzian --- the
Jackiw-Teitelboim gravity action in the discrete setting --- carries
the dynamics.
Independently, the partition function selects flat geometry at the
saddle point, exhibits emergent spatial locality in the two-point
correlator, and concentrates fluctuations at the causal boundary.
In $d = 4$ on the grid $[2]^4$ (3{,}938 convex subsets), the
fluctuation spectrum has a 2-fold degenerate dominant eigenmode;
this degeneracy is explained by the $S_4$ symmetry of the hypercube
and cannot be attributed to graviton physics without data on larger
grids. The thermal entropy fits the JT prediction $S(\beta) =
a/\sqrt{\beta} + S_0$ with $R^2 = 0.99$ at $n = 15$, exhibiting
a smooth crossover from a high-entropy phase to flat spacetime with
no Hawking-Page transition, as expected in 2D.
The density of states is consistent with Hagedorn-like truncation: a Gaussian envelope
in log-space ($R^2 = 0.996$) fits better than the JT
$\sinh\sqrt{E}$ prediction ($R^2 = 0.991$), with a hard UV cutoff
at the Dilworth width of the poset.
A preliminary test in $d = 3$ shows $R^2 = 0.51$ for the $R$-only
fit at small $n$; the falsifiable prediction is that this value
\emph{persists} at larger $n$ in $d \geq 3$ (where EH is dynamical),
unlike $d = 2$ where it collapses to $< 0.01$.
\end{abstract}

\maketitle

%==================================================================
\section{Introduction}
%==================================================================

A companion paper~\cite{DiFioreBH} derives the Bekenstein-Hawking
area law and Hawking temperature from a purely combinatorial
computation: the number of causally convex subsets of the
$d$-dimensional grid $[m]^d$ grows as $\exp(\Theta(m^{d-1}))$
(the dimension law, proved for all $d \geq 2$
in~\cite{DiFioreGrid}), and restricting to time-symmetric
(cylindrical) configurations reduces the exponent by one, yielding
area scaling. That derivation is kinematic: it counts states without
weighting them by an action.

The present paper introduces dynamics. We define a partition function
\begin{equation}
\label{eq:Z}
Z(\beta) = \sum_{S \in \mathrm{CC}(P)} e^{-\beta \cdot S_{\mathrm{BD}}[S]},
\end{equation}
where $\mathrm{CC}(P)$ is the set of causally convex subsets of a
finite poset $P$, and $S_{\mathrm{BD}}[S]$ is the Benincasa-Dowker
action~\cite{BenincasaDowker} restricted to the subregion $S$.
The sum is finite, every term is computable, and the state space
$\mathrm{CC}(P)$ consists of causally consistent subregions --- not
arbitrary subsets. This is a partition function over causally convex subregions of a fixed background poset.

In two dimensions, the Einstein-Hilbert action $\int R \sqrt{g}\,
d^2x$ is topological (proportional to the Euler characteristic) and
carries no local dynamics. The dynamical content of 2D gravity
resides in the conformal mode and the boundary, as captured by the
Jackiw-Teitelboim (JT) action~\cite{Jackiw1985,Teitelboim1983}
and related models~\cite{Saad2019}. A natural
question is whether the partition function~(\ref{eq:Z}) is consistent
with this structure.

We find evidence that it is. On 2{,}061 variable-width grids at
$n = 15$ and $n = 20$, the effective action extracted from $Z(\beta)$ decomposes
into a topological (Einstein-Hilbert) term and a conformal-mode
kinetic (Schwarzian) term, with $R^2 \approx 0.75$ and both coefficients
scaling linearly with $\beta$. The conformal-mode term dominates the
topological term by a factor of $3.1$, as expected in 2D where the
Einstein-Hilbert action is non-dynamical.


%==================================================================
\section{Setup}
%==================================================================

\subsection{Variable-width grids}

A \emph{variable-width grid} with height $T$ and width profile
$(w_0, w_1, \ldots, w_{T-1})$ is the poset
\[
P = \{(i, t) : 0 \leq t < T,\; 0 \leq i < w_t\}
\]
with the componentwise partial order. When all $w_t = m$, this
reduces to the standard grid $[m] \times [T]$.

A subset $S \subseteq P$ is \emph{causally convex} if whenever
$a, b \in S$ with $a \leq c \leq b$, then $c \in S$. Let
$\mathrm{CC}(P)$ denote the set of all causally convex subsets
of $P$.

\subsection{The Benincasa-Dowker action}

The Benincasa-Dowker action~\cite{BenincasaDowker} on a causal set
$S$ is
\begin{equation}
\label{eq:BD}
S_{\mathrm{BD}}[S] = |S| - |\{(x,y) \in S \times S : x \lessdot y\}|,
\end{equation}
where $x \lessdot y$ denotes a link (covering relation): $x < y$
with no $z$ satisfying $x < z < y$.
Equivalently, $S_{\mathrm{BD}}[S] = \chi_{\mathrm{graph}}(\mathrm{Hasse}(S))$, the
graph Euler characteristic of the Hasse diagram: vertices minus edges
(formally verified~\cite{leanrepo}: \texttt{bd\_eq\_euler}).
This is the graph $\chi$, not the cell-complex $\chi$; the
difference counts the 2-cells (``area'').
On a poset embedded in a
Lorentzian manifold, the continuum limit of $S_{\mathrm{BD}}$
is consistent with the Einstein-Hilbert action $\int R \sqrt{g}\,
d^dx$~\cite{BenincasaDowker,Surya2019}.
The positive energy theorem (proved in Lean~4~\cite{leanrepo}:
\texttt{bd\_action\_ge}, \texttt{bd\_action\_eq\_iff\_full}) states that the
full grid uniquely minimizes $\chi_{\mathrm{graph}}$ among all nonempty convex
sub-posets. The action admits a local curvature decomposition:
defining $\kappa(x) = 2 - \deg(x, S)$, the handshaking lemma gives
$2\, S_{\mathrm{BD}} = \sum_{x \in S} \kappa(x)$
(proved in Lean~4~\cite{leanrepo}: \texttt{discrete\_gauss\_bonnet}).
On $[m]^2$: corners have $\kappa = 0$, edge points $\kappa = -1$,
interior points $\kappa = -2$. Flat space minimizes total curvature
because all interior vertices have maximal degree.

\subsection{Discrete curvature and the Schwarzian}

For a variable-width grid with profile $(w_t)$, we define:
\begin{itemize}
\item The \emph{discrete curvature}
$R = \sum_{t} (w_{t+1} - 2w_t + w_{t-1})$, the discrete Laplacian
of the width profile.
\item The \emph{log-Schwarzian}
$\mathrm{LogSch} = \sum_t (\log w_{t+1} - \log w_t)^2$, the
discrete kinetic energy of the conformal mode $\phi_t = \log w_t$.
\end{itemize}

On a 2D manifold with metric $ds^2 = dt^2 + w(t)^2 dx^2$, the
conformal factor is $\phi = \log w$. The kinetic energy
$\int (\partial_t \phi)^2 dt$ is the Schwarzian / JT action in
conformal gauge. The log-Schwarzian is its discrete analogue.


%==================================================================
\section{Saddle Point and Spatial Locality}
%==================================================================

\subsection{Flat geometry dominates}

On the standard grid $[m]^d$, the Benincasa-Dowker action
$S_{\mathrm{BD}}$ is minimized at the full grid (the ``flat''
configuration):
\begin{equation}
S_{\mathrm{BD}}([m]^d) = -(d{-}1)m^d + d\,m^{d-1}.
\end{equation}
For $d = 2$, this simplifies to $-m^2 + 2m = -(m{-}1)^2 + 1$.
This is the unique minimum (proved in Lean~4~\cite{leanrepo}:
\texttt{bd\_action\_ge}, \texttt{bd\_action\_eq\_iff\_full}).
For $\beta > 0$, the partition
function~(\ref{eq:Z}) is dominated by this configuration: flat
spacetime is the saddle point.

On $[4]^2$: the full grid has $S_{\mathrm{BD}} = -8$
(\texttt{bd\_action\_4x4}, formally verified). The anti-diagonal antichain
$\{(0,3),(1,2),(2,1),(3,0)\}$ has $S_{\mathrm{BD}} = +4$
(\texttt{antidiag\_bd\_4}, formally verified --- this also
disproves the earlier claim that $S_{\mathrm{BD,max}} = \min(d,m)$).
The density of states peaks at $S_{\mathrm{BD}} \approx 1$;
the flat configuration is exponentially rare but exponentially
favored by the Boltzmann weight.

\subsection{Emergent spatial locality}

The occupation field $\phi(x) = \mathbf{1}[x \in S]$ defines a
connected two-point function
\begin{equation}
G(x,y) = \langle \phi(x)\phi(y) \rangle - \langle \phi(x) \rangle \langle \phi(y) \rangle,
\end{equation}
where $\langle \cdot \rangle$ denotes the thermal average with
respect to $Z(\beta)$ at $\beta = 1$.

\begin{table}[h]
\caption{Connected correlator $G(r)$ averaged over pairs at
Manhattan distance $r$, at $\beta = 1$.}
\label{tab:correlator}
\begin{ruledtabular}
\begin{tabular}{ccc}
Distance $r$ & $G$ ($m=3$) & $G$ ($m=4$) \\
\midrule
0 & 0.202 & 0.117 \\
1 & 0.090 & 0.048 \\
2 & 0.042 & 0.021 \\
3 & 0.022 & 0.010 \\
4 & 0.017 & 0.006 \\
\end{tabular}
\end{ruledtabular}
\end{table}

The correlator decays monotonically with distance
(Table~\ref{tab:correlator}). No metric, Laplacian, or
nearest-neighbor coupling was introduced --- the decay is a
consequence of the convexity constraint, which correlates the
inclusion of nearby elements more strongly than distant ones.

\subsection{Boundary dominance}\label{sec:boundary}

The occupation probability $\langle \phi(x) \rangle$ is
approximately $0.96$ for interior elements and $0.48$ for corners
on $[m]^2$, and $0.999$ for interior elements and $0.921$ for
corners on $[2]^4$. Fluctuations concentrate at the causal
boundary. In the $d = 2$ setting, only the two extreme corners
(global minimum and maximum of the partial order) can be removed
from the full grid while preserving convexity, yielding exactly 2
boundary degrees of freedom --- matching the boundary-particle
content of JT gravity (proved in Lean~4~\cite{leanrepo}:
\texttt{boundary\_fluctuation}).


%==================================================================
\section{The Controlled Curvature Test}
\label{sec:curvature}
%==================================================================

To disentangle the dependence of $\log Z$ on geometric invariants,
we generated variable-width grid profiles with fixed element count.
An initial sample of 42 profiles at $n = 15$ was extended to
1{,}092 profiles at $n = 15$ and 969 profiles at $n = 20$
(2{,}061 total, computed across 120 cores in 66 seconds).
Profiles were constructed by systematic enumeration of
width sequences $(w_0, \ldots, w_{T-1})$ with $\sum w_t = n$,
$T$ ranging from 3 to $n/2$, and $w_t \geq 1$.

For each profile, $|\mathrm{CC}(P)|$ and $Z(\beta)$ were computed
by exact enumeration. Code and data are available at~\cite{jtrepo}.

\subsection{Einstein-Hilbert alone is insufficient}

A linear fit of $\log Z(\beta{=}1)$ against $R$ alone yields
$R^2 = 0.27$. The sign is correct: more negative $R$ (flatter
geometry) produces larger $\log Z$ (more entropy). But two profiles
with the same $R = -3$ give $\log Z = 6.52$ and $8.41$ --- a factor
of $e^2$ difference at identical integrated curvature. The full width
profile matters, not just its integral.

In retrospect, this is expected: in two dimensions, the
Einstein-Hilbert action $\int R \sqrt{g}\, d^2x$ is topological
(the Gauss-Bonnet theorem). It cannot capture local dynamics. The
dynamical content resides in the conformal mode.

\subsection{Saddle-point decomposition}

The partition function admits a natural decomposition into the
saddle-point action and the entropy of fluctuations:
\begin{equation}
\label{eq:decomp}
\log Z = -\beta \cdot S_{\mathrm{BD}}(\text{full}) + \log\!\left(\frac{Z}{e^{-\beta \cdot S_{\mathrm{BD}}(\text{full})}}\right).
\end{equation}
The first term is the action of the dominant configuration (the full
grid); the second is the entropy of all fluctuations around it.

For a variable-width grid with profile $(w_t)$, the full-grid BD
action is $S_{\mathrm{BD}}(\text{full}) = n - [(n-T) + \sum
\min(w_t, w_{t+1})]$, which depends primarily on~$R$.

Regressing each term separately against geometric invariants on the
full dataset reveals an orthogonal decomposition:

\begin{table}[h]
\caption{Saddle-point decomposition of $\log Z$ ($\beta = 1$).
Original: 42 profiles at $n = 15$. Massive: 1{,}092 profiles at
$n = 15$, 969 at $n = 20$.}
\label{tab:decomp}
\begin{ruledtabular}
\begin{tabular}{lcccc}
Component & Predictor & $R^2$ (42) & $R^2$ (1092) & $R^2$ (969) \\
\midrule
Action & $R$ & 0.19 & 0.10 & 0.12 \\
Entropy & LogSch & 0.69 & 0.54 & 0.62 \\
\midrule
Combined & $R{+}$LogSch & 0.85 & 0.71 & 0.76 \\
\end{tabular}
\end{ruledtabular}
\end{table}

The original 42-profile sample gave $R^2 = 0.85$; the massive
dataset gives $R^2 \approx 0.71$--$0.76$ at $\beta = 1$, rising to
$\approx 0.80$--$0.82$ at $\beta = 3$. The effect is real
(statistically robust on $> 1000$ profiles) but weaker than the
original sample suggested. The log-Schwarzian dominates
($R^2 \approx 0.58$ alone); the discrete curvature $R$ contributes
only $\sim 10\%$.

This is a regression model with a physical interpretation: the curvature $R$ controls the
saddle-point action; the log-Schwarzian controls the entropy of the
fluctuation spectrum. The near-orthogonality ($R^2 = 0.10$ for the
cross-term) means the two effects are essentially independent.

A natural initial hypothesis --- that the log-Schwarzian controls the
entropy by determining the number of near-full configurations
(those with $|S| \geq n - 3$) --- was tested and falsified: the
near-full count is uncorrelated with $\mathrm{LogSch}$
($R^2 < 0.001$). Instead, the log-Schwarzian controls the total
entropy across \emph{all} scales of fluctuation. Profiles with
rapidly varying widths have a richer fluctuation spectrum at every
scale, not just at the boundary.

In two dimensions, the Einstein-Hilbert action
$\int R \sqrt{g}\, d^2x$ is topological, so the action term ($R$)
captures only the topological contribution to $\log Z$. The dynamical
content --- the conformal mode kinetic energy --- resides entirely
in the entropy term ($\mathrm{LogSch}$). This is the structure
of Jackiw-Teitelboim gravity:
\begin{equation}
\label{eq:fit}
\log Z(\beta) = a(\beta) \cdot R + b(\beta) \cdot \mathrm{LogSch} + c(\beta).
\end{equation}

\begin{table}[h]
\caption{Regression of $\log Z$ against geometric invariants.}
\label{tab:regression}
\begin{ruledtabular}
\begin{tabular}{lcccc}
Model & $n$ & $R^2$ & Params & $\beta$-linear \\
\midrule
$R$ only & 15 & 0.27 & 1 & --- \\
$R + \mathrm{LogSch}$ & 15 & 0.85 & 2 & $>0.99$ \\
\midrule
$R$ only & 30 & 0.002 & 1 & --- \\
$\mathrm{LogSch}$ only & 30 & 0.65 & 1 & $>0.99$ \\
$R + \mathrm{LogSch}$ & 30 & 0.65 & 2 & $>0.99$ \\
\end{tabular}
\end{ruledtabular}
\end{table}

At $n = 30$ (10-row grids, 50 profiles), the curvature $R$ alone
explains less than $1\%$ of the variance in $\log Z$, while the
log-Schwarzian explains $65\%$. Adding $R$ to the log-Schwarzian
provides no improvement. Both coefficients retain linear
$\beta$-scaling ($R^2 > 0.99$).

This is the expected 2D behavior. The Einstein-Hilbert action
$\int R \sqrt{g}\, d^2x$ is topological (the Gauss-Bonnet theorem);
its contribution to $\log Z$ does not grow with system size and is
eventually overwhelmed by statistical noise. The $R^2 = 0.27$
observed at $n = 15$ reflects finite-size correlations between $R$
and other geometric invariants that wash out at larger $n$. The
JT/Schwarzian term, by contrast, is dynamical: its coefficient
scales linearly with $\beta$ at both system sizes, and it captures
the dominant variance at $n = 30$.

The remaining $35\%$ of unexplained variance at $n = 30$ is
attributable to higher-order geometric invariants of the width
profile not captured by the two-parameter model. On the original
42 profiles, adding boundary
terms and higher-order invariants improved $R^2$ from
$0.85$ to $0.89$ (Table~\ref{tab:regression15}); the massive
dataset is expected to show a similar relative improvement.

\begin{table}[h]
\caption{Extended regression at $n = 15$ (original 42 profiles).}
\label{tab:regression15}
\begin{ruledtabular}
\begin{tabular}{lccc}
Model & $R^2$ & Params & Interpretation \\
\midrule
$R$ only & 0.27 & 1 & EH (topological) \\
$R + \mathrm{LogSch}$ & 0.85 & 2 & JT action \\
$R + \mathrm{LogSch} + \partial$ & 0.88 & 3 & JT + boundary \\
All invariants & 0.89 & 5 & Diminishing returns \\
\end{tabular}
\end{ruledtabular}
\end{table}


%==================================================================
\section{$\beta$-Scaling and the Effective Action}
\label{sec:beta}
%==================================================================

If $\log Z(\beta)$ genuinely decomposes as
$-\beta \cdot S_{\mathrm{eff}} + S_{\mathrm{entropy}}$, then the
coefficients $a(\beta)$ and $b(\beta)$ in Eq.~(\ref{eq:fit}) must
scale linearly with $\beta$. We fit the model at
$\beta = 0.25, 0.5, 1.0, 2.0, 3.0$:

\begin{table}[h]
\caption{$\beta$-dependence of fit coefficients.}
\label{tab:beta}
\begin{ruledtabular}
\begin{tabular}{ccccc}
$\beta$ & $a(\beta)$ & $b(\beta)$ & $R^2$ & $a/b$ \\
\midrule
0.25 & $-0.04$ & $-0.05$ & 0.17 & 0.78 \\
0.50 & $-0.05$ & $-0.12$ & 0.51 & 0.42 \\
1.00 & $-0.12$ & $-0.30$ & 0.85 & 0.39 \\
2.00 & $-0.27$ & $-0.78$ & 0.89 & 0.35 \\
3.00 & $-0.43$ & $-1.30$ & 0.90 & 0.33 \\
\end{tabular}
\end{ruledtabular}
\end{table}

Linear fits of $a(\beta)$ and $b(\beta)$ against $\beta$ yield:
\begin{align}
a(\beta) &= -0.165\,\beta + 0.035, \quad R^2 = 0.993, \\
b(\beta) &= -0.517\,\beta + 0.169, \quad R^2 = 0.992.
\end{align}
Both coefficients scale linearly with $\beta$ to within $1\%$. The
effective action is therefore
\begin{equation}
\label{eq:Seff}
\boxed{S_{\mathrm{eff}} = 0.165 \cdot R + 0.517 \cdot \mathrm{LogSch}.}
\end{equation}

The intercepts ($0.035$ and $0.169$) are the entropic contributions
at $\beta = 0$. The fit quality increases with $\beta$ (from $0.17$
at $\beta = 0.25$ to $0.90$ at $\beta = 3$): at lower temperature,
the action dominates over entropy fluctuations, and the JT
decomposition becomes sharper.

The ratio $a_0/b_0 = 0.165/0.517 \approx 0.32$ at $n = 15$ sets the
relative strength of the topological and dynamical terms. At $n = 30$,
the topological coefficient shrinks further while the dynamical
coefficient retains its $\beta$-linear scaling ($R^2 > 0.99$),
consistent with the expectation that in two dimensions the
Einstein-Hilbert action is non-dynamical and the physics resides in
the conformal sector.


%==================================================================
\section{Density of States and 2D Black Hole Thermodynamics}
\label{sec:dos}
%==================================================================

\subsection{Density of states}

The density of states $\rho(E) = |\{S \in \mathrm{CC}(P) :
S_{\mathrm{BD}}[S] = E\}|$ on $[4]^2$ is fit by several models in
log-space:

\begin{table}[h]
\caption{Density of states fits on $[4]^2$ (log-space).}
\label{tab:dos}
\begin{ruledtabular}
\begin{tabular}{lcc}
Model & $R^2$ & Interpretation \\
\midrule
Gaussian $\log\rho = ax^2{+}bx{+}c$ & 0.996 & Hagedorn \\
JT $\rho \sim \sinh\sqrt{E}$ & 0.991 & JT gravity \\
Exponential & 0.962 & --- \\
Power law & 0.957 & --- \\
\end{tabular}
\end{ruledtabular}
\end{table}

The JT prediction $\rho \sim \sinh(2\pi\sqrt{2\gamma E})$ is a good
approximation ($R^2 = 0.991$), but the best fit is a Gaussian
envelope ($R^2 = 0.996$): exponential growth in $\sqrt{E}$
terminated by a hard UV wall. This is the density of states of a
consistent with Hagedorn-like behavior --- exponential growth capped by a maximum
energy --- rather than pure JT gravity.

The UV cutoff has a combinatorial origin: the maximum BD action on
any convex subset of $[m]^2$ is $m$, achieved uniquely by the
anti-diagonal antichain (proved in Lean~4~\cite{leanrepo}).
For $[2]^4$, the observed maximum is $w = 6$
(the middle layer $\binom{4}{2}$); the general
$S_{\mathrm{BD,max}} = w$ for $[m]^d$ with $d \geq 3$ is
conjectured but not yet proved. This is a hard boundary in
the spectrum --- a UV completion built into the combinatorial
structure.

The BD energy spectrum on $[m]^2$ is now bracketed by proved
theorems at both ends:

\textbf{Theorem (Spectral gap).}
For any nonempty convex $S \subseteq [m]^2$ with
$S_{\mathrm{BD}}(S) \neq S_{\mathrm{BD}}(\text{flat})$,
we have $S_{\mathrm{BD}}(S) \geq S_{\mathrm{BD}}(\text{flat}) + 1$.
The gap $\Delta = 1$ is universal for all~$m$.
The first excited level has degeneracy exactly~$2$
(remove the global minimum or maximum).
Proved in Lean~4~\cite{leanrepo}.

\textbf{Theorem (Top state uniqueness).}
The maximum BD action on $[m]^2$ is $S_{\mathrm{BD,max}} = m$,
achieved by exactly one convex subset: the anti-diagonal antichain.
Proved in Lean~4~\cite{leanrepo}.

\noindent
The spectral gap is the discrete analogue of the mass gap in
JT gravity: the ground state (flat spacetime) is separated from
all excitations by a universal energy quantum. The two boundary
degrees of freedom in the first excited level match the
JT boundary particle content (spectral gap theorem +
boundary fluctuation theorem of \S\ref{sec:boundary}).

\subsection{2D black hole thermodynamics}

The thermal entropy $S(\beta) = \log Z + \beta \langle E \rangle$
is well-fit by the JT gravity prediction $S(\beta) = a/\sqrt{\beta}
+ S_0$:

\begin{table}[h]
\caption{JT entropy fit $S(\beta) = a/\sqrt{\beta} + S_0$.}
\label{tab:entropy}
\begin{ruledtabular}
\begin{tabular}{lccc}
Grid & $a$ & $S_0$ & $R^2$ \\
\midrule
$[2] \times 5$ ($n=10$) & 5.92 & $-2.42$ & 0.946 \\
$[3] \times 5$ ($n=15$) & 7.36 & $-3.72$ & 0.987 \\
\end{tabular}
\end{ruledtabular}
\end{table}

The fit quality increases with system size ($R^2 = 0.95 \to 0.99$),
as expected for a continuum limit. The JT coupling $\gamma$,
extracted from the coefficient $a$, grows with grid size
($\gamma = 0.44$ at $n = 10$, $\gamma = 0.69$ at $n = 15$),
reflecting the increasing number of degrees of freedom.

The specific heat peaks at $C \approx 4.4$ at $\beta \approx 1.1$
(for $n = 15$), with the crossover temperature scaling as
$T_c \sim 1/S_0$: the product $\beta_c \cdot \log|\mathrm{CC}|$ is
approximately constant ($\approx 7$--$8$) across grid sizes. This is
a smooth crossover, not a discontinuous transition --- consistent
with JT gravity, which has no Hawking-Page transition in
2D~\cite{Saad2019}.


%==================================================================
\section{The $d = 4$ Fluctuation Spectrum}
\label{sec:d4}
%==================================================================

On the grid $[2]^4$ (16 elements, $|\mathrm{CC}| = 3{,}938$ convex
subsets), the correlation matrix $G(x,y)$ at $\beta = 1$ has the
following eigenvalue structure:

\begin{table}[h]
\caption{Eigenvalues of the correlation matrix on $[2]^4$.}
\label{tab:eigenvalues}
\begin{ruledtabular}
\begin{tabular}{ccc}
Eigenvalue & Degeneracy & $S_4$ irrep \\
\midrule
0.080 & 2 & 2-dim (weight-2 sector) \\
0.010 & 6 & standard $\times$ 2 \\
0.001 & 6 & standard $\times$ 2 \\
\end{tabular}
\end{ruledtabular}
\end{table}

The correlation matrix commutes with the $S_4$ permutation group of
the hypercube (maximum commutator $\sim 10^{-16}$). The 2-fold
degenerate dominant eigenmode is the 2-dimensional irreducible
representation of $S_4$ acting on the weight-2 sector, not a
graviton signal. Any $S_4$-invariant matrix on 16 vertices will
exhibit this structure.

To test whether graviton-like excitations emerge beyond the symmetry
prediction, the computation must be extended to $[3]^4$ (81
elements), where the $S_4$ representation decomposes differently.
If the 2-fold dominance persists with a gap that grows with $m$,
that would be evidence for graviton physics. If the dominant
degeneracy changes to match the new $S_4$ decomposition, the
$[2]^4$ result was purely representation-theoretic.

The correlator on $[2]^4$ decays as $G(r) \sim 0.003/r^{2.8}$.
The exponent exceeds the $d = 4$ graviton prediction $\alpha = 2$,
but with only 5 distinct distance values on a 2-element grid,
finite-size effects preclude a meaningful comparison.


%==================================================================
\section{Discussion}
%==================================================================

The partition function $Z(\beta) = \sum_{S \in \mathrm{CC}(P)}
e^{-\beta \cdot S_{\mathrm{BD}}[S]}$ is a well-defined discrete
partition function over causally convex subregions. Its state space is finite,
its action is the Benincasa-Dowker functional, and the
sum is exact. No continuum limit, regularization, or gauge-fixing
is required.

In two dimensions, the effective action extracted from $Z(\beta)$
is consistent with the structure of Jackiw-Teitelboim gravity. At $n = 15$,
both $R$ and the log-Schwarzian contribute, with a saddle-point
decomposition (Table~\ref{tab:decomp}) showing the two predictors
are nearly orthogonal (cross-$R^2 = 0.10$). At $n = 30$, the
topological term $R$ washes out ($R^2 < 0.01$) while the
log-Schwarzian retains $R^2 = 0.65$ with $\beta$-stable coefficients
($R^2 > 0.99$). This is the correct asymptotic behavior for 2D:
the Einstein-Hilbert action is topological and its finite-size
contribution shrinks, leaving the JT/Schwarzian as the sole
dynamical term.

This result is specific to $d = 2$, where the Einstein-Hilbert
action is topological and the dynamics resides in the conformal
mode. A preliminary test in $d = 3$ on fixed-$N$ variable-width
grids gives $R^2 \approx 0.51$ for the $R$-only fit at small $n$.
The critical test is whether this $R^2$ \emph{persists} at larger
$n$ in $d = 3$ (where EH is dynamical), in contrast to $d = 2$
where it collapses. If so, the partition function correctly
distinguishes topological from dynamical gravity across dimensions.

The five established results --- saddle-point selection of flat
geometry, emergent spatial locality, boundary dominance of
fluctuations, the JT action decomposition with $\beta$-stable
coefficients, and the persistence of the Schwarzian term at larger
system sizes --- are independent of the $S_4$ eigenvalue analysis
on $[2]^4$, which remains inconclusive pending larger grids.

The BD-weighted CC partition function provides a concrete, finite,
computable discrete statistical model with gravity-like features.
It makes no reference to a background metric, a continuum manifold,
or quantum field theory. The objects summed over are the causally
convex subregions of a finite poset; the action is the
Benincasa-Dowker functional; the result, in $d = 2$, is
consistent with Jackiw-Teitelboim gravity.

\textbf{Limitations.}
A massive-scale test (2{,}061 profiles across $n = 15$ and $n = 20$)
confirmed the JT-like decomposition but at lower $R^2$ than the
original 42-profile sample ($R^2 \approx 0.75$ vs $0.85$ at
$\beta = 1$). The original sample was somewhat favorable; the
effect is real but modest. The JT identification is
regression-based, not derived from the partition function analytically.
The density-of-states comparison (Gaussian $R^2 = 0.996$ vs JT
$R^2 = 0.991$) is marginal on a tiny system ($[4]^2$, 1146 subsets)
and should not be taken as definitive evidence for universality.
Scaling to $[m]^2$ with $m \geq 6$ and to $[m]^3$ is needed to
test whether the observed structure persists or is a finite-size
artifact.

\textbf{Open questions.} (i)~\emph{Dimensional dependence of $R^2$.}
In $d = 2$, the $R$-only variance explained decreases with $n$
($0.27$ at $n = 15$, $< 0.01$ at $n = 30$), consistent with EH
being topological. In $d \geq 3$, EH is dynamical: the
$R$-only $R^2$ should \emph{stabilize} at a nonzero value as $n$
grows. A preliminary test at $d = 3$ gives $R^2 = 0.51$ at small
$n$; confirming that this does not collapse at larger $n$ would
demonstrate that the partition function correctly distinguishes
topological from dynamical gravity.
(ii)~\emph{SYK correspondence.} JT gravity is dual to the SYK
model~\cite{SYK}. The spectral form factor of $Z(t) = \sum
e^{-(\beta + it) S_{\mathrm{BD}}}$ on $[4]^2$ exhibits a Gaussian
slope ($R^2 = 0.997$) and qualitative slope-dip-ramp structure, but
the level statistics are crystalline (not Wigner-Dyson) due to the
small number of distinct energy levels. Ensemble-averaging over grid
profiles or larger grids may be needed to access the chaotic regime.
(iii)~Computation of $|\mathrm{CC}([3]^4)|$ to resolve the
graviton-vs-$S_4$ question.
(iv)~The Lorentzian partition function
$Z = \sum e^{i S_{\mathrm{BD}}}$ requires a contour prescription,
following Sorkin's quantum measure approach~\cite{Sorkin1994}.


\begin{acknowledgments}
The combinatorial framework builds on~\cite{DiFioreCAG,DiFioreGrid}.
The formal verification of the underlying dimension law
uses the Mathlib library~\cite{mathlib}; source code at~\cite{leanrepo}.
\end{acknowledgments}


\begin{thebibliography}{99}

\bibitem{DiFioreCAG}
T.~DiFiore,
Causal-algebraic geometry: A Grothendieck-type framework for partial
orders, preprint (2026).

\bibitem{DiFioreGrid}
T.~DiFiore,
Order-convex subsets of grid posets: A new exponential combinatorial
class, preprint (2026).

\bibitem{DiFioreBH}
T.~DiFiore,
Black hole thermodynamics from counting convex subsets of a lattice,
preprint (2026).

\bibitem{BenincasaDowker}
D.~M.~T.~Benincasa and F.~Dowker,
Phys.\ Rev.\ Lett.\ \textbf{104}, 181301 (2010).

\bibitem{Surya2019}
S.~Surya,
Living Rev.\ Relativ.\ \textbf{22}, 5 (2019).

\bibitem{Jackiw1985}
R.~Jackiw, in \emph{Quantum Theory of Gravity},
edited by S.~M.~Christensen (Hilger, Bristol, 1984), pp.~403--420.

\bibitem{Teitelboim1983}
C.~Teitelboim,
Phys.\ Lett.\ B \textbf{126}, 41 (1983).

\bibitem{Saad2019}
P.~Saad, S.~H.~Shenker, and D.~Stanford,
arXiv:1903.11115 [hep-th] (2019).

\bibitem{Sorkin1994}
R.~D.~Sorkin,
Mod.\ Phys.\ Lett.\ A \textbf{9}, 3119 (1994).

\bibitem{SYK}
A.~Kitaev,
talk at KITP (2015);
J.~Maldacena and D.~Stanford,
Phys.\ Rev.\ D \textbf{94}, 106002 (2016).

\bibitem{mathlib}
The Mathlib Community,
\emph{Mathlib: The Lean Mathematical Library},
\url{https://github.com/leanprover-community/mathlib4} (2024).

\bibitem{leanrepo}
T.~DiFiore,
\emph{Causal-Algebraic Geometry: Formal Proofs},
\url{https://github.com/tomdif/causal-algebraic-geometry-lean} (2026).

\bibitem{jtrepo}
T.~DiFiore,
\emph{JT Gravity from Convex Subsets: Code and Data},
\url{https://github.com/tomdif/jt-gravity-from-convex-subsets} (2026).

\end{thebibliography}

\end{document}
